% Owned by the polygons workstream.  Result IDs P1--P8 are stable.
\section{A.3. The image of the polygon invariant}
\label{sec:A3}

We keep the arrangement labelled.  Let $E=\{1,\ldots,n\}$, let
$\widetilde E=\{+i,-i:i\in E\}$, and assume that $\Sigma$ is a finite,
non-pencil arrangement of distinct tame pseudolines in $\mathbb {RP}^{2}$.
Multiple vertices are allowed.  The associated oriented matroid $M$ has
rank three and is simple; ``simple arrangement'' would impose the stronger
condition that $M$ is uniform.  We distinguish these two uses of
``simple.''  We choose orientations of the old pseudolines and use a
chirotope representative $\chi_M:E^3\longrightarrow\{-1,0,1\}$
calibrated with these signed half-line labels.  For example, for oriented
projective lines represented by normal vectors $a_i$, the calibration is
$\chi_M(i,j,k)=\operatorname{sgn}\det(a_i,a_j,a_k)$, and in an affine
chart the positive endpoint of the line with normal $(u_i,v_i,w_i)$ has
direction $(-v_i,u_i)$.  The geometric proofs below give the same
calibration for pseudolines.  Changing this convention requires changing
the corresponding signed labels as well; it is not an additional
ambiguity of the relative position.

An unoriented relative position of a new pseudoline $d$, distinct from
every old line, is a simple rank-three extension $N$ on $E\sqcup\{d\}$
with $N\setminus d=M$, modulo reorientation of $d$ only.  We write
$q(i,j)=\chi_N(i,j,d)$.  Thus $q$ is the chirotope of $N/d$, and
reorientation of $d$ changes $q$ to $-q$.  These correspondences use the
rank-three topological representation theorem: an oriented matroid
determines the labelled cell decomposition, and any two representations
with the same data are related by a homeomorphism.  We checked the
projective formulation in \cite[\S7.1.4, pp.~6--7]{RGZpreprint}, and the
relative uniqueness statement in \cite[Theorems~4--5, pp.~15,18--20]{BKMS2002}.
The latter fixes the old arrangement cellwise.  Consequently a
representation obtained abstractly for $N$ can be carried to the given
$\Sigma$; we do not replace the prescribed arrangement by a different
realization without this comparison.

We make explicit the projective comparison used here.  On the
spherical double covers, equal covector data identify the labelled
vertices, edges, and two-dimensional cells.  Their closures are
balls, by \cite[Theorem~2, pp.~10--11]{BKMS2002}.  Match the vertices,
extend over each pair of antipodal edges by interval homeomorphisms,
and then over each pair of antipodal two-cells by disk
homeomorphisms extending the already matched boundary circles.
On the second member of each pair use the antipodal conjugate of
the first map.  No cell is its own antipode, since its nonzero sign
vector differs from its negative.  The resulting map is an
antipodal-equivariant cellular homeomorphism of spheres, and hence
descends to a labelled homeomorphism of the projective arrangements.
This is the comparison carrying a deletion representation to the
prescribed $\Sigma$ in every use below.

\subsection{An explicit finite image criterion}

A signed boundary word is a cyclic listing $W$ of $\widetilde E$, each
label occurring once.  We may temporarily orient the cyclic listing;
reversing it is allowed when the deployed polygon is unoriented.  From
$W$ define an alternating map $q_W:E^2\longrightarrow\{-1,0,1\}$.
For $i\ne j$, put $q_W(i,j)=0$ if the two endpoint pairs
$\{+i,-i\}$ and $\{+j,-j\}$ do not alternate.  If they alternate, put
$q_W(i,j)=+1$ when $+j$ lies in the positively directed open arc from
$+i$ to $-i$, and put $q_W(i,j)=-1$ otherwise.  This is an entirely
combinatorial rule; in particular, it does not assume that $W$ is valid.

For any alternating $q:E^2\longrightarrow\{-1,0,1\}$ define
$\chi_q:(E\sqcup\{d\})^3\longrightarrow\{-1,0,1\}$ by
\begin{equation}
 \chi_q(i,j,k)=\chi_M(i,j,k),\qquad
 \chi_q(i,j,d)=q(i,j),                                      \label{eq:P-combined}
\end{equation}
and extend by alternation, with value zero on repeated entries.  The
following version of the chirotope test is explicit also in the
nonuniform case.  Its nonzero unordered triples must be the bases of a
rank-three matroid: for bases $B,C$ and $b\in B\setminus C$, some
$c\in C\setminus B$ must make $(B\setminus\{b\})\cup\{c\}$ a basis.
For every five distinct labels $\lambda,a,b,c,e$, the three signs
\begin{equation}
 \chi_q(\lambda,a,b)\chi_q(\lambda,c,e),\quad
 -\chi_q(\lambda,a,c)\chi_q(\lambda,b,e),\quad
 \chi_q(\lambda,a,e)\chi_q(\lambda,b,c)                    \label{eq:P-GP}
\end{equation}
must either include both signs or all be zero.  These are the two
conditions (CHI1)--(CHI2) in the opened source
\cite[\S7.2.3, p.~10]{RGZpreprint}.  The unsigned basis condition is
essential when zeros occur; it must not silently be omitted.

We next specify the cluster expansion used in the criterion.  Suppose
that $\chi_q$ passes this test and no row of $q$ is zero.  The rank-two
matroid of $q$ has no loops.  Its signed elements have a cyclic order of
parallel-ray classes, with opposite classes halfway round.  Choose one
of the two cyclic orientations.  In a signed ray class containing
$\varepsilon i$ and $\eta j$, choose $k$ outside the corresponding
unoriented parallel class.  Order the two labels by
\begin{equation}
 \varepsilon i \text{ precedes }\eta j
 \quad\Longleftrightarrow\quad
 -\eta\,\chi_M(i,j,k)q(i,k)=+1.                         \label{eq:P-expansion}
\end{equation}
Place the opposite class in the reversed, sign-negated order.  Denote
the resulting word by $\operatorname{Exp}_M(q)$.  The proof of P3 below
shows that the relation is a total order inside each packet and is
independent of $k$.  This gives a practical construction without real
coordinates: choose a signed representative for each class in a
half-circle of the rank-two order, sort the classes using the signs of
$q$, and sort each packet by \eqref{eq:P-expansion}.

\begin{resultenv}{P4}{Proved}{Complete image criterion}
A cyclic signed word $W$ is the deployed boundary polygon of a
non-supercritical relative position for the prescribed $\Sigma$ if and
only if the following finite tests hold, with the calibration just
specified:
\begin{enumerate}
 \item $\chi_{q_W}$ satisfies the unsigned basis axiom and
       \eqref{eq:P-GP};
 \item no row of $q_W$ is identically zero;
 \item $W$ agrees, up to a cyclic shift, with
       $\operatorname{Exp}_M(q_W)$.
\end{enumerate}
If $W$ is given without a cyclic orientation, perform these steps with
either orientation; the resulting extension is taken modulo
reorientation of $d$.  The tests include arbitrary vertex
multiplicities and any number of vertices on $d$.
\end{resultenv}

\begin{proof}
For an actual relative position, cut along a narrow M\"obius band about
$d$.  Two remaining chords cross exactly when their intersections with
$d$ were distinct.  By the Jordan separation property for a proper arc
of a disk, they cross exactly when their endpoints alternate.  In that
case their signed cyclic order gives $\chi_N(i,j,d)$, by the
rank-two chirotope convention.  Thus the pair rule gives $q_W=q$.
The combined map is the chirotope of $N$, and it passes the two axioms.
A zero row at $i$ would say that $\{i,d\}$ has rank one: $d$ would be
parallel to $i$ in the oriented matroid, and hence the same projective
pseudoline.  This is excluded.  Finally the local disk computation in
P3 is exactly the expansion of the packets, so the third condition
holds.

Conversely, the first condition produces a rank-three oriented matroid
$N$ with chirotope $\chi_{q_W}$, whose deletion of $d$ is exactly $M$.
The second condition makes $d$ a nonloop and prevents it from being
parallel to any old element.  All old elements were already distinct
and nonparallel.  Therefore the projective topological representation
of $N$ consists of the prescribed number of distinct pseudolines.
Its deletion represents $M$, so the relative uniqueness/homeomorphism
comparison stated above carries that deletion to $\Sigma$ and carries
the new pseudoline to a relative position for $\Sigma$.  P3 computes
its boundary word as $\operatorname{Exp}_M(q_W)$; the third condition
identifies that word with $W$.  This proves sufficiency and exhausts
the image.  Reversal changes the contraction orientation, and no old
label is reoriented in this comparison.
\end{proof}

There is an equivalent localization formulation that explains precisely
where Las Vergnas's theorem enters.  For independent $i,j\in E$ let
$C_{ij}$ be the signed cocircuit with entries
$C_{ij}(k)=\chi_M(i,j,k)$.  Given $q$, set
\begin{equation}
               \sigma(C_{ij})=q(i,j).                    \label{eq:P-sigma}
\end{equation}
First require that this is well defined when two ordered pairs define
the same cocircuit, and that $\sigma(-C)=-\sigma(C)$.
For every rank-one flat of $M$ contract that flat.  On the cyclic list
of $2m$ cocircuits of the resulting rank-two oriented matroid, the
values of $\sigma$ must, up to cyclic shift, be one of
\[
 0^{2m},\qquad +^{m-1}0-^{m-1}0,\qquad +^{m}-^{m}.
\]
These tests are exactly Las Vergnas's localization criterion.  We
opened its full version, including zero values, in
\cite[Definition~6.4 and Theorem~6.5, p.~25]{GP2021}; that source states
the dual version for colocalizations.  To apply it here we take $M^*$.
A subset $A\subseteq E$ has nullity two in $M^*$ precisely when
$E\setminus A$ has rank one in $M$, since
$r_{M^*}(A)=|A|-3+r_M(E\setminus A)$.  Moreover
$(M^*|A)^*=M/(E\setminus A)$, and circuits of the first matroid are
cocircuits of the second.  Thus its restrictions are exactly the
rank-two contractions above, with the same cyclic sign patterns.
This proves that the criterion applies in the direction required here,
including the nonuniform cases.  The nonzero localization determines
the extension; P2 supplies a direct uniqueness proof from the
chirotope.  The image criterion P4 adds the explicit cluster expansion
to recover the literal deployed polygon, rather than merely its
contraction.

\subsection{The proposed three local conditions are insufficient}

The phrase ``three triangles'' needs a test that can be applied before
the existence of a relative position is known.  Otherwise it already
assumes the conclusion.  We use the strongest natural local test on a
stable candidate word: an adjacent signed pair
$(\varepsilon i,\eta j)$ is a candidate empty triangle at $v_{ij}$ if,
for every $k\notin\{i,j\}$, all nonzero signs in
\begin{equation}
 \chi_M(i,j,k)q(i,j),\quad \varepsilon q(i,k),\quad
 \eta q(j,k)                                               \label{eq:P-triangle}
\end{equation}
agree.  In an actual simple arrangement these are the signs of the
vertex and the two boundary endpoints relative to the $k$th chord.
They agree exactly when that chord does not cut off the alleged
triangle.  Adjacency excludes an endpoint in its base, and testing all
other chords excludes a crossing of either side.  Hence
\eqref{eq:P-triangle} is an exact empty-triangle test for a valid word.
We apply this same finite necessary test to candidates that may fail
global compatibility.

\begin{resultenv}{P5}{Refuted, with exact replacement P4}{Lead L3}
The $2n$ weight rule, the cluster-reversal rule, and the presence of at
least three empty-triangle candidates, even tested against every old
line by \eqref{eq:P-triangle}, do not characterize the image.  This
fails for a simple arrangement of six straight lines.  The complete
replacement is P4, or equivalently the localization criterion together
with the prescribed packet expansion.
\end{resultenv}

\begin{proof}
Let $a_i=(1,t_i,t_i^2)$ for $t_i=i-1$, $1\le i\le6$, and let the old
projective line $i$ have equation $a_i\cdot x=0$.  For $i<j<k$,
\[
 \det(a_i,a_j,a_k)=(t_j-t_i)(t_k-t_i)(t_k-t_j)>0.
\]
Thus all old triples have sign $+1$, all old lines are distinct, and
no three old lines are concurrent.  In particular the arrangement is
nontrivial.  Consider the cyclic word
\[
 W=(1,2,-3,-4,-6,-5,-1,-2,3,4,6,5).
\]
It is stable and exactly antipodal: the second half is the negative of
the first.  Its rank-two pair signs, in lexicographic order
$12,13,14,15,16,23,24,25,26,34,35,36,45,46,56$, are
\[
 (+,-,-,-,-,-,-,-,-,+,+,+,+,+,-).
\]
The three adjacent arcs
\[
                (2,-3),\qquad(-1,-2),\qquad(3,4)
\]
pass \eqref{eq:P-triangle} for every other old line.  For the first
arc all four sign triples are $(-,-,-)$; for the other two arcs all
four sign triples are $(+,+,+)$.  The remaining nine adjacent arcs
fail, so there are exactly three candidates.  Give each marked arc
weight one.  In their cyclic order the three complementary arcs have
weights $4,1,4$, the numbers of boundary edges they contain.  The
total weight is $1+1+1+4+1+4=12=2n$.  There are no nontrivial clusters,
so the cluster rule holds vacuously.

Nevertheless, in the combined map of \eqref{eq:P-combined}, take
$\lambda=5$, $a=1$, $b=3$, $c=6$, and $e=d$.  The three terms in
\eqref{eq:P-GP} are all $-1$.  For example they are
\[
 \chi_M(5,1,3)q(5,6)=-1,\quad
 -\chi_M(5,1,6)q(5,3)=-1,\quad
 q(5,1)\chi_M(5,3,6)=-1.
\]
This violates the chirotope axiom, so no pseudoline in the prescribed
arrangement has the candidate word.  The exact-integer program
\texttt{polygons\_l3\_search.py} checks all determinants, all twelve
triangle tests, the stated violation, and a failure of the full signed
basis-exchange axiom.  Its complete output is
\texttt{audit/polygons\_l3\_search\_output.json}.  A second direct
verifier, \texttt{polygons\_l3\_verify.py}, repeats the certificate
without searching.  Thus the counterexample concerns the prescribed
$M$, not a polygon borrowed from a different arrangement.
\end{proof}

The search also tested every stable signed word for the four- and
five-element moment-curve arrangements (48 and 384 oriented words),
and all 3840 such words for the six-line test arrangement of A.1;
none of those searches produced this obstruction.  The six-element
moment-curve example was found at candidate 49.  These finite negative
searches do not assert minimality among all arrangements.

\subsection{Recovering the weighted graph}

\begin{resultenv}{P6}{Proved}{Reconstruction of vertices, arrows, rotations, and weights}
For a fixed labelled nontrivial arrangement $\Sigma$, its full image of
deployed polygons determines all non-supercritical relative positions,
all ordinary specialization arrows, their directions and weights, and
their cyclic incidences at each position.  Adjoining the known $n$
supercritical positions recovers the corridor arrows and their
weights.  Together with the prescribed transition transport used in
constructing the surface, these data determine its cellularly embedded
weighted graph, up to label-preserving isomorphism of cellular maps.
The reconstructed graph alone does not assert that every face has
order four; that is the separate local-link assertion of A.1.
\end{resultenv}

\begin{proof}
P2 below reconstructs the extension at every polygon.  Its signed
cocircuit localization is \eqref{eq:P-sigma}.  A projective old vertex
$v$ is a pair $\{C_v,-C_v\}$ of cocircuits; write
$F_v=C_v^0\subset E$ for the set of old lines through it, and
$m_v=|F_v|$.  An ordinary immediate specialization is obtained by
replacing the two values at $\pm C_v$ by zero and leaving every other
localization value unchanged, provided the resulting localization
passes the extension criterion and is not supercritical.  This rule
is finite and orientation-independent: simultaneous negation of all
values just reorients $d$.

We verify the geometric equivalence of the rule.  An elementary sweep
to $v$ moves no other old vertex across $d$, so it has exactly this
effect on the localization.  Conversely, if the rule yields an
extension $N'$, represent it with its new line through $v$.  Choose a
disk about $v$ containing no other old vertex.  Resolving the new line
to either side inside that disk gives the two localizations obtained
by replacing the zero at $C_v$ by the two possible signs; all other
coordinates are unchanged.  Exactly one has the given nonzero sign
of the original extension $N$.  Uniqueness of the extension therefore
identifies that local resolution with $N$.  This comparison is a
homeomorphism preserving the old arrangement cellwise, so the local
resolution proves that the required empty fan exists for the original
relative position.  In a cut-open disk it is precisely the
specialization triangle, subdivided by the $m_v$ concurrent old
lines.  There are $m_v-1$ boundary edges in its base, which proves
that the arrow weight is $m_v-1$.

At $N'$, the class $F_v$ is visible directly from the endpoint
nonalternation relation.  Its two boundary packets give the two
generization intervals.  Reversing one packet performs the
corresponding resolution; the other packet stays fixed.  At $N$, the
same marked interval is the specialization base.  Thus the signed
boundary word specifies the interval occupied by each incident
ordinary arrow, not only its two end vertices.  Their circular order
is the required local rotation, and the lengths of the complementary
intervals give all corner weights.  The intervals partition the
$2n$ boundary edges between transition bases and corner intervals,
which also proves the $2n$ weight identity.

For a supercritical position $i$, use the parallel extension
$N_i$ with $d$ parallel to $i$.  Its localization is
$\sigma_i(C)=C(i)$, and its contraction entries are
$q_i(j,k)=\chi_M(j,k,i)$.  Let $N$ be substable, with its unique old
vertex $v$ on line $i$.  There is a corridor arrow $N\to N_i$ exactly
when, after choosing the sign of the new element,
\[
          \sigma_i(C)\ne0 \quad\Longrightarrow\quad
          \sigma_N(C)=\sigma_i(C) .                     \label{eq:P-corridor}
\]
Necessity is immediate by collapsing the corridor.  For sufficiency,
contract $i$.  The cyclic rank-two localization criterion shows that,
with the zero at $v$ prescribed, there are just the two choices
corresponding to the two local sides of $i$ at $v$.  A narrow strip
about $i$, with the new line crossing $i$ at $v$ and following either
side elsewhere, realizes these choices.  Its localization away from
$i$ is $C(i)$; on $i$ it is the prescribed rank-two pattern.
Therefore \eqref{eq:P-corridor} and uniqueness identify $N$ with one
of these narrow-strip positions.  This is the required corridor.
The $n-m_v$ lines not through $v$ cross the corridor as disjoint
transversal chords, dividing its side into $n-m_v+1$ edges.  Hence
the ordinary corridor weight is $n-m_v+1$; its modified weight at
the supercritical vertex is $m_v-1$, and their sum is $n$.
The order of the corridor flags at $i$ is the order of the old
vertices on the twofold cover of $i$, also determined by $M/i$.

Finally, a graph with cyclic orders at its vertices and specified
identifications of the two local orientation choices across its edges
is a signed rotation system.  Thicken each vertex to a disk and each
edge to a band with that identification.  Its boundary cycles are
determined by following the finite permutations of flags; attaching
one disk to each such cycle gives the unique cellular map with these
data.  Transition transport is part of the construction being
reconstructed, and its flag identifications are obtained from the
labelled local sweep or the corridor.  This establishes the asserted
reconstruction without deducing quadrangularity or embeddability of
the original large disks from the mere existence of the signed
rotation system.
\end{proof}

\section{A.4. Proofs of the disk criterion and the two polygon assertions}
\label{sec:A4}

\subsection{The precise disk criterion}

\begin{resultenv}{P1}{Proved, with the one-class qualification}{Sufficiency in R1}
Let $\Delta$ be a closed disk, with finitely many labelled tame proper
arcs $\delta_i$ whose $2n$ endpoints on $\partial\Delta$ are distinct.
Suppose each pair meets at most once, at an interior transverse
crossing, and disjointness together with equality is an equivalence
relation on the arcs.  If there are at least two equivalence classes,
this system is the cut-open system of a projective pseudoline
arrangement and an added pseudoline.  If there is just one class,
the same assertion holds exactly when its endpoints form two
consecutive packets, each containing one endpoint of each chord.
The reconstructed old arrangement is nontrivial exactly when its
pairwise intersection points are not all the same point.
\end{resultenv}

\begin{proof}
We first prove that the packet condition is automatic with at least
two classes.  Fix a class $C$ and an arc $\beta$ in another class.
Every arc in $C$ crosses $\beta$ once.  The pairwise disjoint arcs of
$C$ therefore separate the two endpoints of $\beta$ and occur in a
linear chain along it.  To see explicitly that there is no branch,
cut $\Delta$ along $\beta$.  In either resulting half-disk the pieces
of arcs in $C$ are pairwise disjoint arcs joining distinct points of
the cut to the opposite boundary; noncrossing forces the same linear
order, reversed when the boundary is traversed in the opposite
direction.  Thus the endpoints of $C$ occur in two packets with
opposite orders.

No endpoint of an arc outside $C$ can lie strictly between two
consecutive endpoints in either packet.  Indeed those two arcs of
$C$ bound a middle strip.  An arc starting in that strip and crossing
each of its two boundary arcs exactly once cannot end in a component
outside both: after the first crossing it is outside one boundary,
and to reach the other boundary it would have to recross the first.
This contradicts the one-crossing hypothesis.  Since every outside
arc crosses both boundaries, neither endpoint can start in the
middle strip.  It follows that the two packets are consecutive among
all chord endpoints, not merely among endpoints belonging to $C$.

Collapse each packet to a point of the boundary.  The quotient of the
disk is again a disk.  For completeness, one may describe the quotient
map in polar coordinates: choose a nondecreasing degree-one boundary
map $\phi$ that is constant exactly on the packet intervals, and use
$t\mapsto (1-r)t+r\phi(t)$ on the circle of radius $r<1$.  These maps
are strictly increasing for $r<1$, extend continuously to the boundary
collapse, and induce a continuous bijection from the quotient onto a
disk.  Compactness gives a homeomorphism.  Finite piecewise linear
choices give the corresponding tame local model.

Choose a representative chord from each of the $r$ classes.  They
cross pairwise.  Their $2r$ boundary points have order
$C_1,\ldots,C_r,C_1,\ldots,C_r$: after cutting at the two endpoints
of $C_1$, every other chord has one endpoint in either half, and two
such chords cross exactly when their orders in the two halves agree.
This proves the displayed order.  Consequently the involution pairing
the two points of each class is an order-preserving, fixed-point-free
boundary involution.  Identify this boundary by that involution.
After a boundary homeomorphism it is the usual antipodal
identification, so the quotient is $\mathbb {RP}^2$.

The image of the boundary is a new pseudoline $d$.  Each old chord
closes to a simple curve, with its two endpoints identified and no
other self-identification.  It crosses $d$ once, hence is an essential
one-sided simple curve and is a pseudoline.  Chords from different
classes retain their unique interior transverse intersection.  Chords
in the same class meet only at the boundary vertex created for that
class.  Their opposite packet orders give, in the local disk after
gluing, the cyclic order of rays of a pencil: each pair of old curves
crosses there, and each crosses $d$.  Thus all required pairwise
intersection and crossing conditions hold.  Expanding the packet
vertices inside a narrow band about $d$ recovers the original disk
system, up to a labelled homeomorphism.

For one class the same construction works when two packets exist.
Conversely, if every old chord is disjoint, every pair of old
pseudolines must have met on $d$.  Since each old pseudoline meets
$d$ only once, all those intersections are a single vertex, and a
small neighborhood of that vertex gives exactly two packets.  The
condition is necessary.  Finally the quotient construction lists all
old intersection points, interior crossings and class vertices;
their being a single common point is precisely the pencil case.
\end{proof}

The packet condition occurs explicitly in the transcription of
154/035, after the disjointness-equivalence clause.  The transcription
marks some bracketed wording about antipodality and uniqueness as
doubtful or deleted.  P1 does not depend on either reading of that
bracket: it proves the packet assertion where it is automatic, and
states its exact necessity in the remaining case.  If R1 is read as
asserting sufficiency of the two abbreviated conditions for all finite
systems without this qualification, it is false: three disjoint small
chords cutting off three separate boundary caps have word
$(1,-1,2,-2,3,-3)$, and admit no two packets.  The finite endpoint
certificate is in \texttt{polygons\_certificates.py}.  This degenerate
example does not arise from the nontrivial rank-three arrangements in
the question.  For $n=0$ the empty disk system may be closed with an
arbitrary added line; for $n=1,2$ every old arrangement is a pencil.
Those cases are outside the nontrivial rank-three hypothesis.  No
assertion about infinite or nontame systems is being used.

\subsection{Injectivity and cluster reversal}

\begin{resultenv}{P2}{Proved}{R2: the literal polygon determines the relative position}
For the fixed labelled nontrivial arrangement $\Sigma$, the map from
non-supercritical relative positions to literal deployed signed
polygons is injective.  The polygon determines the rank-two
contraction by the signed pair-alternation rule, including all its
parallel classes.  The only ambiguity is reorientation of the added
element, corresponding to reversal of the chosen cyclic orientation.
\end{resultenv}

\begin{proof}
The Jordan separation argument in P4 recovers every value
$q(i,j)=\chi_N(i,j,d)$ from the word.  In particular, no guess at a
cluster decomposition is needed: $i$ and $j$ are in the same cluster
if and only if their endpoint pairs do not alternate.  For an actual
word this is exactly equality of their intersections with $d$, hence
is an equivalence relation.  We know $\chi_N$ on every old triple
from $M$, and on every triple containing $d$ from $q$.  Alternation
and zeros on repetitions determine all remaining ordered triples.
Thus \eqref{eq:P-combined} determines the entire extension.  The
relative topological representation comparison identifies its
relative position in the fixed arrangement.  Reversing $q$ changes
exactly the chirotope entries containing $d$, which is precisely
reorientation of $d$.

Equivalently, $q$ gives the cocircuit localization by
\eqref{eq:P-sigma}, and Las Vergnas's uniqueness of a single-element
extension with a prescribed localization gives the same conclusion.
The direct argument above identifies the exact data involved and
does not need to invoke the stronger existence theorem.  In
particular it avoids confusing the contraction with the deployed
polygon: the latter has the additional within-packet orders dealt
with in P3.
\end{proof}

\begin{resultenv}{P3}{Proved}{R3: the exact antipodal permutation}
If the vertices of $\Sigma$ on $d$ have line clusters of sizes
$k_1,\ldots,k_\nu$, the deployed word consists of paired consecutive
packets in opposite order.  Let $H$ be the half-turn on the $2n$
positions and let $R$ reverse the order separately within every
packet, fixing singleton packets.  The label-negation involution is
$H\circ R=R\circ H$.  It equals $H$ exactly when $d$ is stable.
The within-packet order is \eqref{eq:P-expansion}.
\end{resultenv}

\begin{proof}
At an old vertex $v$ on $d$, choose a small disk containing no other
old vertex.  The $k$ old line germs and the germ of $d$ form a pencil
of transverse proper arcs.  A homeomorphism of this disk carries
the finite star to straight radial segments, preserving its cyclic
order.  Take $d$ locally to be $y=0$ and write the other lines as
$x=b_i y$, with distinct $b_i$.  The two nearby boundary arcs meet
them at $x=\epsilon b_i$ and $x=-\epsilon b_i$.  Their orders are
reversed.  This proof is topological: the straightening is only of
one local star, and is available whether the whole arrangement is
stretchable or not.

After the packets are collapsed, the endpoints are the rank-two
cyclic order on $d$.  Corresponding packets are opposite and have
the same cardinality.  Hence the full word has the form
\[
       B_1\cdots B_r\,(-B_1)^{\rm rev}\cdots(-B_r)^{\rm rev}.
\]
Its first and second halves each have $n$ entries.  The formula for
label negation follows entry by entry.  If a packet has at least two
entries, its reversal is nonidentity, so label negation is not the
half-turn.  If every packet is a singleton, the reversals are all
identities.  Thus the equivalence with stability is exact.

We also verify the signed expansion formula, needed for the
algorithm.  It involves only two lines $i,j$ of one cluster, $d$,
and an old line $k$ outside that cluster.  This pencil-with-transversal
configuration is straightenable by choosing an affine chart with
$d$ at infinity.  Write its normal rows as $a_i=(w_i,z_i)$,
$a_j=(w_j,z_j)$, $a_k=(w_k,z_k)$ with $w_\ell\in\mathbb R^2$.
After reorienting $i,j$ by $\varepsilon,\eta$, their vectors
$\varepsilon w_i$ and $\eta w_j$ lie on the same ray. Write
$w_j=\lambda w_i$, where $\lambda\ne0$ and
$\operatorname{sgn}\lambda=\varepsilon\eta$. Elementary expansion gives
\[
 \det(a_i,a_j,a_k)=-(z_j-\lambda z_i)\det(w_i,w_k).
\]
On a large circle of radius $R$, the signed endpoints are near
$J(\varepsilon w_i)/\|w_i\|=J(\eta w_j)/\|w_j\|$, where $J$
rotates by $\pi/2$. Their angular offsets are ordered by
$\varepsilon z_i/\|w_i\|$ and $\eta z_j/\|w_j\|$.
Consequently the sign of the second offset minus the first is the sign of
$-\eta\chi_M(i,j,k)q(i,k)$, as in
\eqref{eq:P-expansion}.  The order is geometric and independent of
the choice of transversal $k$, so the formula is independent of $k$
and is transitive for the whole packet.  Only the local pencil and
transversal were straightened, which proves the formula for arbitrary
pseudoline arrangements.

A smallest nontrivial explicit failure of the half-turn claim is
given by old equations $x=0$, $x=z$, $y=0$ and $d:z=0$.  The old
normal determinant is $+1$, so the old arrangement is nontrivial.
The word is $(2,1,3,-1,-2,-3)$; its half-turn takes $2$ to $-1$.
The exact computation in \texttt{polygons\_certificates.py} verifies
the arrangement, the word, and the full chirotope axiom.  Therefore
the literal antipodality assertion on 154/018 is false for
nonstable positions; the correction R3 is true.
\end{proof}

\subsection{A local fact used in the surface analysis}

\begin{resultenv}{P7}{Proved}{Commuting resolutions at distinct vertices}
Let $d$ be a non-supercritical relative position through distinct old
vertices $v_1,\ldots,v_\nu$.  Resolving any chosen subset to either
local side can be performed independently.  The resulting
generizations are indexed by $\{0,+,-\}^{\nu}$, where zero means
unresolved, and the $2^\nu$ fully resolved positions are distinct.
In particular two distinct chosen vertices give four distinct
positions in each ordinary commuting square.
\end{resultenv}

\begin{proof}
Distinct vertices on $d$ cannot share an old line: that old line
would meet $d$ twice unless it were $d$, contrary to the
non-supercritical assumption.  Choose mutually disjoint small disks
about the $v_i$, meeting no other old vertices.  In each disk replace
the arc of $d$ by either of the two nearby transverse arcs, or leave
it unchanged.  Outside these disks leave $d$ fixed.  The resulting
curve stays simple and meets each old line once, because an old line
is affected in at most one of the disks and the local replacement
retains its unique crossing.  Operations in disjoint disks commute.
The localization changes at exactly the corresponding pair of old
cocircuits, from zero to the chosen sign.  Different index vectors
therefore give different localizations and different relative
positions, even after forgetting the orientation of $d$: an old
vertex not on $d$ fixes the relative global sign, or, equivalently,
one can use an unchanged nonzero entry of the contraction between two
different clusters.  Such an entry exists because $M$ has rank three
and $d$ is non-supercritical.  This proves the asserted cube and the
distinctness of its square vertices.
\end{proof}

P7 proves the ordinary commuting part of a corner argument, but does
not itself identify every corner of the disk-glued space with such a
square.  In particular two specialization bases can share an old
line, so their future apices cannot both lie on a distinct added
pseudoline; that case must be treated by the corridor/supercritical
construction.  We retain this distinction to avoid using a local
cube as an unproved global gluing theorem.

\begin{resultenv}{P8}{Proved}{The exposed-fan or corridor alternative}
Suppose that a non-supercritical new pseudoline $d$ passes through
exactly one old vertex $v$.  On the spherical double cover, let $\alpha$ be either
open semicircle of $d$ between the two lifts of $v$.  Choose an old
line $i$ through $v$ adjacent to $d$ on one side, and let $L$ be the
closed lune bounded by $\overline\alpha$ and the corresponding
semicircle of $i$.  If $L$ contains an old vertex in its interior,
then an empty specialization fan has its base in $\alpha$.  If it
contains no old interior vertex, $d$ is in the corridor position next
to $i$ on that side.  In particular, a half-circle between its
generization ports contains an
ordinary specialization or is an exceptional corridor half-circle.
\end{resultenv}

\begin{proof}
Here ``adjacent'' refers to the cyclic order of the old line germs
and the germ of $d$ at $v$.  An old line through $v$, other than $i$,
cannot enter the interior of $L$: it has no further intersection
with either boundary line, and its germ is outside this sector.
Every old line not through $v$ meets each wall once and cuts $L$ in
one proper arc joining its $d$ wall to its $i$ wall.  Indeed each
wall semicircle contains exactly one of the two antipodal
intersections with that old line, and these are the only boundary
intersections.  We call these arcs the wires and direct them from
the $d$ wall to the $i$ wall.

We justify that the resulting finite directed graph of old vertices
has no directed cycle.  Let $N$ be the oriented matroid of the whole
arrangement including $d$, use $g=d$ as the distinguished element
at infinity, and use $f=i$ as objective.  Thus the affine space in
the terminology of \cite[pp.~30--31]{SZ1993} is the hemisphere
$S_d^+$ containing $L$.  We choose the sign of $i$ so that the
program orientation points toward the $i$ wall in this lune.  The
definition of its graph $G_i$ assigns a direction to an edge along
a wire by the sign of $i$ at the wire's point at infinity on $d$.
That sign is constant on the chosen open $d$ wall, and is nonzero:
a traversing wire does not pass through $v=d\cap i$.  Hence every
wire edge is directed, and its direction is precisely the one just
specified.  The interior graph in $L$ is a subgraph of $G_i$.

The program satisfies the hypotheses of the cited definition.  In
particular $d$ is not a loop and $i$ is not a coloop.  For the latter,
another old line through $v$, an old line not through $v$, and $d$
form a rank-three basis after deleting $i$; such a line not through
$v$ exists because the old arrangement is not a pencil.  The proof
of \cite[Corollary~4.5, p.~40]{SZ1993} shows that no rank-three
oriented matroid program is non-Euclidean.  By its Definition~3.8(i)
and Proposition~3.9 (p.~33), such a program has no directed cycle
(a path may use undirected edges, but a directed cycle must contain
at least one directed edge).  In our subgraph every edge is directed,
so it is an ordinary finite directed acyclic graph.

If there is an interior vertex, choose a source $s$ of this graph.
For each old line through $s$, its incoming arc from the $d$ wall
has no old vertex before $s$; otherwise the last such vertex and
the intervening edge would contradict the source property.  These
incoming arcs have distinct endpoints on $\alpha$ and are mutually
disjoint except at $s$.  They form a star that partitions $L$ into
regions.  Exactly one region has the $i$ wall on its boundary.  Each
outgoing branch at $s$ must lie in that region: if it initially lay
in any other region, then to reach its endpoint on the $i$ wall it
would have to cross an incoming arc again or meet the $d$ wall
again.  Either event contradicts the unique-intersection property
of the pseudolines.  Thus the incoming branches are consecutive in
the cyclic order at $s$.  Their two extreme branches and the
interval of $\alpha$ between their endpoints bound a fan disk $F$.

No old line not incident to $s$ enters the interior of $F$.  Such a
line cannot cross either extreme incoming arc, since an intersection
there would be a vertex preceding $s$.  If it entered through the
base on $d$, it would have to leave through one of those extreme
arcs or through $d$ a second time; it cannot leave through $s$
because it is not incident to $s$.  These alternatives are all
impossible.  The intermediate incoming arcs only subdivide $F$
into the usual empty triangles.  We can therefore move the base
arc of $d$ through this empty fan to $s$, without changing any
other old intersection.  This gives the required ordinary
specialization and places its full base in $\alpha$.

If there is no old interior vertex, all wires in the open lune are
pairwise disjoint.  Move the $d$ wall to a curve just inside the
$i$ wall, keeping its endpoints at the two lifts of $v$.  This move
can be performed without passing an old vertex: the disjoint
wire arcs cut the lune into finitely many disk strips, with fans
allowed at their endpoints on $i$, and within each such strip we
slide the boundary interval across a rectangle.  The slides agree
on their common wire sides.  Equivalently they give an isotopy of
the new arc through transverse arcs, preserving one intersection
with each wire.  It extends over the complementary cells to a
homeomorphism preserving every labelled old line.  The final arc
is a narrow push-off of $i$ crossing it at $v$, on the specified
local side.  This is exactly the corridor position.  Conversely,
such a corridor lune contains no old interior vertex.  The two
alternatives prove the statement and justify the exceptional case
described on 154/093 independently of its doubtful marginal text.
\end{proof}
